% !TEX root = ../main.tex
% !TEX program = lualatex
\documentclass[../main.tex]{subfiles}
\IfSubfilesClassLoaded{%
  \externaldocument{\subfix{../build/main}}%
}{}
% =============================================================================
% CHAPTER 11: CONTRACTING FUNDAMENTALS
% DESCRIPTION: Binding contract elements, FAR system overview, and KO/PM roles.
% =============================================================================
\begin{document}
\IfSubfilesClassLoaded{\chapter{EDO Study Guide}}{}
%====================
% Contents here
%====================

% === INTRO TO CONTRACTING (EDO 3.2.1) ===
\section{Intro to Contracting Fundamentals}

\subsection{Summary}
\begin{itemize}
  \item Contracts protect the contractor and the Government by establishing a binding legal relationship.
  \item Essential elements are mutual assent, consideration, capacity, and lawful purpose.
  \item Governing documents include the \ac{far}, \ac{dfars}, and agency supplements.
  \item The three \ac{ko} types are \ac{pco}, \ac{aco}, and \ac{tco}.
  \item \acp{pm} manage the program and are accountable to the \ac{mdauthority}; \acp{ko} ensure contract performance and compliance.
  \item CICA implements competition requirements, and sole-source actions require approved justification.
\end{itemize}
~\autocite{edo-3-2-1-intro-contracting-2025}

\subsection{Why Contracts Are Needed}
\begin{itemize}
  \item Establish a legal relationship between the Government and the contractor.
  \item Define the rights and responsibilities of each party.
  \item Allow changes within the terms and conditions of the legal relationship.
\end{itemize}
\note{Contracts protect both the contractor and the Government throughout the acquisition process~\autocite{edo-3-2-1-intro-contracting-2025}.}

\subsection{Essential Elements of a Contract}
\begin{description}
  \item[\textbf{Mutual assent.}] A contract is formed when an offer is accepted by the offeree, creating a meeting of the minds on material terms such as price, quantity, and delivery~\autocite{edo-3-2-1-intro-contracting-2025}.
  \item[\textbf{Consideration.}] Each party provides something of value (e.g., performance, payment, schedule relief); courts consider sufficiency and adequacy in Government contracts~\autocite{edo-3-2-1-intro-contracting-2025}.
  \item[\textbf{Capacity.}] Parties must be legally competent and authorized to bind their organizations; incapacity or lack of authority undermines enforceability~\autocite{edo-3-2-1-intro-contracting-2025}.
  \item[\textbf{Lawful purpose.}] The objective must be legal and enforceable; contracts contrary to public policy or statute are void~\autocite{edo-3-2-1-intro-contracting-2025}.
\end{description}
\note{Without all elements properly satisfied, a contract may not be valid in the eyes of the court~\autocite{edo-3-2-1-intro-contracting-2025}.}

\subsection{Government Relationship to Contractors}
\begin{description}
  \item[\textbf{Transactional.}] Contractors seek profit while the Government demands cost, schedule, and performance results; avoid unauthorized commitments~\autocite{edo-3-2-1-intro-contracting-2025}.
  \item[\textbf{Professional.}] Both sides must understand the contract and work cohesively to meet requirements~\autocite{edo-3-2-1-intro-contracting-2025}.
  \item[\textbf{Collaborative.}] Contractors bring technical expertise, while programmatic decision-making remains inherently governmental~\autocite{edo-3-2-1-intro-contracting-2025}.
  \item[\textbf{Constrained.}] Ethics rules and contractual obligations govern interactions; avoid conflicts of interest~\autocite{edo-3-2-1-intro-contracting-2025}.
\end{description}

\subsection{Federal Acquisition Regulation System}
\begin{itemize}
  \item The system establishes policies and procedures for contracting and procurement across executive agencies~\autocite{edo-3-2-1-intro-contracting-2025}.
  \item It consists of:
  \begin{itemize}
    \item The \ac{far}.
    \item The \ac{dfars} (\ac{dod} \ac{far} Supplement), including \ac{dod} source selection procedures.
    \item Agency supplements and acquisition regulations such as \ac{nmcars}, \ac{som}, the \ac{navsea} Source Selection Guide, and the \ac{navsea} Contracts Handbook.
  \end{itemize}
\end{itemize}
\note{Contracting officers use the \ac{far} and its supplements to carry out the contracting process~\autocite{edo-3-2-1-intro-contracting-2025}.}

\subsection{Types of Contracting Officers}
\begin{itemize}
  \item \ac{pco}: Handles procurement from pre-solicitation through award and signs the contract on behalf of the Government~\autocite{edo-3-2-1-intro-contracting-2025}.
  \item \ac{aco}: Performs contract administration functions as assigned by the \ac{pco} (see \ac{far}~42.302(a))~\autocite{edo-3-2-1-intro-contracting-2025}.
  \item \ac{tco}: After a termination notice, negotiates an equitable settlement with the contractor~\autocite{edo-3-2-1-intro-contracting-2025}.
\end{itemize}
\note{A single \ac{ko} can serve in all three roles depending on organizational needs and contract type~\autocite{edo-3-2-1-intro-contracting-2025}.}

\subsection{KO and PM Roles and Responsibilities}
\begin{description}
  \item[\textbf{KO responsibilities.}] Ensure contract performance and compliance; serve as the contract advisor for the program~\autocite{edo-3-2-1-intro-contracting-2025}.
  \item[\textbf{KO major functions.}] Prepare solicitations and contracts; communicate with offerors and negotiate; ensure consistency with \ac{far}/\ac{dfars}; award, administer, modify, and terminate contracts; legally bind the Government through warranted authority~\autocite{edo-3-2-1-intro-contracting-2025}.
  \item[\textbf{PM responsibilities.}] Manage the program and remain accountable to the \ac{mdauthority} for cost, schedule, and performance in accordance with the \ac{dod} 5000 series~\autocite{edo-3-2-1-intro-contracting-2025}.
\end{description}
\note{The \ac{ko}, \ac{pm}, and their teams must understand each other's responsibilities and governing regulations~\autocite{edo-3-2-1-intro-contracting-2025}.}

\subsection{Differences Between the PM and the KO}
\begin{longtblr}[
    label   = {tab:pm_vs_ko_differences},
    caption = {Differences between the PM and the KO},
    entry   = {PM vs KO},
    remark{\tSource} = {\tabCite{edo-3-2-1-intro-contracting-2025}}
  ]
  {colspec = {@{} lll @{}}}
  \toprule
  {Category} & {PM} & {KO} \\
  \midrule
  Authority & Tenure agreement & Warrant \\
  Responsibility & Entire program & Contract \\
  Background/Training & Technical & Business \\
  Guiding directives & DoW 5000 series & FAR system and contract management \\
  Organization & Program Office (IPT) & Matrix (IPT) \\
  \bottomrule
\end{longtblr}
\note{PM authority is provided by the appropriate \ac{cae} in accordance with \ac{dod} 5000 and \ac{dodi}~5000.66; \acp{hca} provide \acp{ko} authority via a Certificate of Appointment (Standard Form~1402) warrant. \ac{ipt} means Integrated Product Team~\autocite{edo-3-2-1-intro-contracting-2025}.}

\subsection{Who's Responsible?}
\begin{longtblr}[
    label   = {tab:responsibility_ko_pm},
    caption = {Contracting responsibilities by role},
    entry   = {KO and PM Responsibilities},
    remark{\tSource} = {\tabCite{edo-3-2-1-intro-contracting-2025}}
  ]
  {colspec = {@{} ll @{}}}
  \toprule
  {Action} & {Role} \\
  \midrule
  Decide which vendor will be awarded the contract & KO / SSA caveat \\
  Conduct contract negotiations with the selected vendor & KO \\
  Set the desired award date for the contract & PM \\
  Set the maximum dollar amount to be awarded & PM \\
  Monitor contractor performance & PM/KO \\
  Make changes to the contract & KO \\
  \bottomrule
\end{longtblr}
\note{The \ac{ko} awards and signs the contract. In formal source selections, the \ac{ssa} selects best value when appointed; under \ac{far}~15.303, the contracting officer may serve as the \ac{ssa} by default~\autocite{edo-3-2-1-intro-contracting-2025,FAR}.}

\subsection{Advantages of Competition}
\begin{itemize}
  \item Incentivizes lower prices for goods and services.
  \item Encourages innovation in transformational technologies.
  \item Promotes higher quality products and services.
  \item Enables best-value tradeoffs for performance improvements.
  \item Provides opportunities for capable small businesses.
  \item Strengthens the defense industrial base.
\end{itemize}
\note{Competition improves both value and capability outcomes~\autocite{edo-3-2-1-intro-contracting-2025}.}

\subsection{Competition in Contracting Act (CICA)}
\begin{description}
  \item[\textbf{Full and open competition.}] All responsible sources may compete for the requirement~\autocite{edo-3-2-1-intro-contracting-2025}.
  \item[\textbf{Exclusion of sources.}] Allows alternate-source strategies, small-business set-asides, and \ac{sba} Section~8(a) actions~\autocite{edo-3-2-1-intro-contracting-2025}.
  \item[\textbf{Seven authorities.}] The coursebook memory hook is the seven statutory authorities for other-than-full-and-open competition (e.g., sole source); for current DoD use, cite the active \ac{rfo} Part~6 deviation and revised \ac{far}~6.103 text instead of relying on legacy 6.302 numbering~\autocite{edo-3-2-1-intro-contracting-2025,FAR-Overhaul-Part6-2026,DoD-RFO-Deviation-Part6-2026}.
  \item[\textbf{Approval.}] Most other-than-full-and-open competition requires a written \ac{ja}; under the active \ac{rfo} structure, justification content and approval levels sit in revised \ac{far}~6.104 with DoD overlays in \ac{dfars}/\ac{pgi}~206.104, while public availability moved to revised \ac{far}~6.301. Public-interest use remains an agency-head determination with congressional notice~\autocite{edo-3-2-1-intro-contracting-2025,FAR-Overhaul-Part6-2026,DoD-RFO-Deviation-Part6-2026}.
\end{description}

\subsubsection{Seven Other-Than-Full-and-Open Authorities}
\begin{enumerate}
  \item Only one responsible source will satisfy agency requirements.
  \item Unusual and compelling urgency.
  \item Industrial mobilization; engineering, developmental, or research capability; or expert services.
  \item International agreement.
  \item Authorized or required by statute.
  \item National security.
  \item Public interest.
\end{enumerate}
\note{The seven-category board hook is still useful, but the current DoD working text is the active \ac{rfo} Part~6/DFARS~206 deviation. Use revised \ac{far}~6.103 for the authority, revised \ac{far}~6.104 for \ac{ja} content and approval levels, and revised \ac{far}~6.301 for posting/public-availability rules~\autocite{edo-3-2-1-intro-contracting-2025,FAR-Overhaul-Part6-2026,DoD-RFO-Deviation-Part6-2026}.}

\subsection{Contracting Terms}
\begin{description}
  \item[\textbf{Sole source.}] A contract negotiated with only one source; other-than-full-and-open competition requires justification and approval (typically a \ac{ja} using \ac{nmcars} format)~\autocite{edo-3-2-1-intro-contracting-2025}.
  \item[\textbf{Responsiveness.}] For sealed bidding, awards are made without discussion; the bid must conform to all \ac{ifb} requirements~\autocite{edo-3-2-1-intro-contracting-2025}.
  \item[\textbf{Responsibility.}] A prospective contractor must be determined responsible, including financial health, to receive an award (applies to \ac{ifb} and \ac{rfp})~\autocite{edo-3-2-1-intro-contracting-2025}.
\end{description}
\note{Tip to remember: you are responsible for your own financial health~\autocite{edo-3-2-1-intro-contracting-2025}.}

\subsection{Determining Responsibility}
The prospective contractor (offeror) must:
\begin{itemize}
  \item Have adequate financial resources or the ability to obtain them.
  \item Comply with the required or proposed delivery schedule.
  \item Maintain a satisfactory performance record.
  \item Maintain a satisfactory record of integrity and business ethics.
  \item Possess (or obtain) the necessary organization, experience, accounting and operational controls, and technical skills.
  \item Have the necessary equipment or the ability to obtain it.
  \item Be otherwise qualified and eligible under applicable laws and regulations.
\end{itemize}
~\autocite{edo-3-2-1-intro-contracting-2025}

\subsection{Determining Responsiveness}
\begin{itemize}
  \item A bid must comply in all material respects with the \ac{ifb} to be considered for award.
  \item Example: if the offered delivery date (including mailing or transmittal time) is later than required, the bid is nonresponsive and rejected.
  \item Example: if the solicitation requires copper but the bidder offers tin, the bid is nonresponsive.
\end{itemize}
~\autocite{edo-3-2-1-intro-contracting-2025}

\subsection{Quick Review}
\begin{itemize}
  \item A contract is valid only when the parties have mutual assent, consideration, capacity, and lawful purpose.
  \item The \ac{ko} binds the Government through a warrant; the \ac{pm} owns program execution but cannot make or modify contractual commitments.
  \item \ac{pco}, \ac{aco}, and \ac{tco} are contracting-officer roles across award, administration, and termination.
  \item Use \ac{far}, \ac{dfars}, and Navy supplements as the regulation stack; policy preference or program urgency cannot replace contracting authority.
  \item For competition questions, start with full and open competition, then explain other-than-full-and-open authority, written \ac{ja}, approval level, and public-interest notice if applicable.
  \item Responsibility is about the offeror's capability and eligibility; responsiveness is about whether a sealed bid conforms to the \ac{ifb}.
\end{itemize}

%====================
% End of file
%====================
\IfSubfilesClassLoaded{
  \RenewDocumentCommand{\entryname}{}{\textbf{\color{Modern} Acronym}}
  \RenewDocumentCommand{\descriptionname}{}{\textbf{\color{Modern} Definition}}
  \printnoidxglossary[
    type=\acronymtype,
    title=Acronyms,
    style=long-booktabs
  ]}{}
\end{document}
